% Pseudocode, math and plain-language explanations of the Bellman-consistency hallucination detector.
% Self-contained: upload this single file to Overleaf and compile with pdfLaTeX.
% Companion to docs/problem_formulation.tex (same notation) and the Markdown files in docs/pseudocode/.
\documentclass[11pt]{article}
\usepackage[margin=1in]{geometry}
\usepackage{amsmath,amssymb,amsthm,booktabs,enumitem,array}
\usepackage{algorithm}
\usepackage{algpseudocode}
\usepackage[most]{tcolorbox}
\usepackage[hidelinks]{hyperref}

\newtheorem{definition}{Definition}
\newtheorem{proposition}{Proposition}
\newcommand{\R}{\mathbb{R}}
\newcommand{\abs}[1]{\left|#1\right|}
\newcommand{\Qbar}{\bar Q}
\newcommand{\sd}{\operatorname{std}}
\newcommand{\med}{\operatorname{median}}
\newcommand{\code}[1]{\texttt{#1}}
\algrenewcommand\algorithmicrequire{\textbf{Input:}}
\algrenewcommand\algorithmicensure{\textbf{Output:}}

\newtcolorbox{inwords}{colback=blue!3,colframe=blue!40!black,boxrule=0.4pt,arc=2pt,
  left=6pt,right=6pt,top=4pt,bottom=4pt,fonttitle=\bfseries\small,title=In words}
\newtcolorbox{example}[1][]{colback=green!3,colframe=green!35!black,boxrule=0.4pt,arc=2pt,
  left=6pt,right=6pt,top=4pt,bottom=4pt,fonttitle=\bfseries\small,title=Small example,#1}

\title{The Bellman-Consistency Audit:\\Pseudocode, Formulas and Explanations}
\author{}
\date{}

\begin{document}
\maketitle

\noindent
This document explains, step by step, what the code computes: every formula, an algorithm box for every procedure,
and a plain-language explanation with a small example. Each part names the function in the repository it
describes. Notation follows \code{docs/problem\_formulation.tex}.

\paragraph{Fixed settings.} Imagined horizon $H=12$ for detection ($24$ and $40$ in the compounding study), audited
steps $L=H-1$, critic heads $K=5$, discount $\gamma=0.95$ (ManiSkill; $0.99$ on DeepMind Control), UCB weight
$\beta=1$.

\tableofcontents
\bigskip

\begin{algorithm}[h]
\caption{End-to-end, for one trained robot}\label{alg:e2e}
\begin{algorithmic}[1]
\State $\theta\gets$ trained TD-MPC2 agent \Comment{frozen; the detector never changes it}
\State $\mathcal{D}\gets$ 50 real episodes recorded with the agent's own planner \Comment{observations, actions, real rewards}
\State $\mathcal{C},\mathcal{T}\gets$ first 25 episodes (calibration), last 25 (test)
\State $b\gets\textsc{CriticBias}(\theta,\mathcal{C})$ \Comment{Sec.~\ref{sec:bias}}
\State $(X_{\mathcal C},\text{rollouts})\gets\textsc{AuditFeatures}(\theta,\mathcal{C},b)$ \Comment{Secs.~\ref{sec:signals},\ref{sec:features}}
\State $(y_{\mathcal C},\varepsilon_{\text{hall}},\varepsilon_{\text{clean}})\gets\textsc{AnswerKey}(\text{rollouts},\mathcal{C})$ \Comment{Sec.~\ref{sec:labels}}
\State $f\gets\textsc{FitTrees}(X_{\mathcal C},y_{\mathcal C})$ \Comment{Sec.~\ref{sec:trees}}
\State $(X_{\mathcal T},\text{rollouts})\gets\textsc{AuditFeatures}(\theta,\mathcal{T},b)$
\State $y_{\mathcal T}\gets\textsc{AnswerKey}(\text{rollouts},\mathcal{T};\ \varepsilon_{\text{hall}},\varepsilon_{\text{clean}})$ \Comment{thresholds from calibration only}
\State $p_{\mathcal T}\gets f(X_{\mathcal T})$ \Comment{hallucination probability per imagined step}
\State \Return $\textsc{StratifiedAUROC}(p_{\mathcal T},y_{\mathcal T})$ \Comment{Sec.~\ref{sec:eval}; ours $\approx0.82$}
\end{algorithmic}
\end{algorithm}

\begin{inwords}
We take a trained robot and never change it. We record 50 real episodes. On the first 25 we learn what ``normal''
looks like and fit the detector. On the other 25 we test it: for every piece of imagined future we know the real
outcome, so we can check whether the detector raised its alarm on the right pieces.
\end{inwords}

%==================================================================================================================
\section{World model and imagined rollout}\label{sec:wm}
%==================================================================================================================

\subsection{The TD-MPC2 parts we use (trained, frozen)}

\begin{center}
\begin{tabular}{@{}lll@{}}
\toprule
Part & Formula & Meaning\\
\midrule
Encoder & $z=h_\theta(o)$ & real observation $\to$ latent vector (there is no decoder)\\
Dynamics & $\hat z'=d_\theta(z,a)$ & deterministic next latent\\
Reward head & $\hat r=R_\theta(z,a)$ & predicted reward (two-hot over 101 bins $\to$ scalar)\\
Critic heads & $Q^{(i)}_\theta(z,a),\ i=1..K$ & $K=5$ value estimates (two-hot $\to$ scalar)\\
Policy prior & $\mu_\pi(z)$ & mean action of the policy\\
Mean critic & $\Qbar(z,a)=\frac1K\sum_i Q^{(i)}_\theta(z,a)$ & \\
Value & $V_\theta(z)=\frac1K\sum_i Q^{(i)}_\theta\bigl(z,\mu_\pi(z)\bigr)$ & value of a state, used everywhere below\\
\bottomrule
\end{tabular}
\end{center}

\paragraph{How TD-MPC2 trains its critic} (\code{third\_party/tdmpc2/tdmpc2/tdmpc2.py::\_update}).
A training sample is a short real slice $o_0..o_3$, $a_0..a_2$, $r_0..r_2$ from the replay buffer:
\begin{align*}
z_0&=h_\theta(o_0)\ \text{(real)},\qquad \hat z_{t+1}=d_\theta(\hat z_t,a_t)\ \text{(imagined)},\\
y_t&=r_t+\gamma\min_{\text{2 random target heads}}Q^{(i)}_{\bar\theta}\bigl(h_\theta(o_{t+1}),\pi(h_\theta(o_{t+1}))\bigr)
\quad\text{(real reward, real next observation)},\\
\mathcal{L}_Q&=\sum_t\rho^t\sum_i\operatorname{CE}\bigl(Q^{(i)}_\theta([z_0,\hat z_1,\hat z_2]_t,a_t),\ y_t\bigr).
\end{align*}
The critic's \emph{inputs} are partly imagined but its \emph{targets} come from real rewards and real next
observations. Our \emph{grounded} variant feeds $h_\theta(o_t)$ instead of $\hat z_t$ as the critic input, with no
gradient into the encoder or dynamics, so the critic is anchored to real states.

\subsection{Imagined rollout (\code{src/auditor/signals.py::audit\_rollout})}

Given a real observation $o_t$ and actions $\mathbf a=(a_0,\dots,a_{H-1})$:
\begin{align*}
\hat z_0&=h_\theta(o_t)\quad\text{(real)},\qquad \hat z_{k+1}=d_\theta(\hat z_k,a_k)\quad\text{(imagined)},\\
\hat r_k&=R_\theta(\hat z_k,a_k),\qquad q^{(i)}_k=Q^{(i)}_\theta(\hat z_k,a_k),\qquad q_k=\tfrac1K\textstyle\sum_iq^{(i)}_k,\\
v^{(i)}_k&=Q^{(i)}_\theta\bigl(\hat z_k,\mu_\pi(\hat z_k)\bigr),\qquad v_k=\tfrac1K\textstyle\sum_iv^{(i)}_k=V_\theta(\hat z_k),
\qquad k=0,\dots,H \ (\text{$\hat r,q$ up to }H-1).
\end{align*}

\begin{algorithm}[h]
\caption{\textsc{Imagine}$(o_t,\mathbf a)$}\label{alg:imagine}
\begin{algorithmic}[1]
\Require real observation $o_t$; actions $a_0..a_{H-1}$
\State $\hat z_0\gets h_\theta(o_t)$ \Comment{the only real point of the rollout}
\For{$k=0$ to $H-1$}
  \State $\hat z_{k+1}\gets d_\theta(\hat z_k,a_k)$
  \State $\hat r_k\gets R_\theta(\hat z_k,a_k)$;\quad $q^{(i)}_k\gets Q^{(i)}_\theta(\hat z_k,a_k)$ for $i=1..K$;\quad $q_k\gets\operatorname{mean}_i q^{(i)}_k$
\EndFor
\For{$k=0$ to $H$}
  \State $v^{(i)}_k\gets Q^{(i)}_\theta(\hat z_k,\mu_\pi(\hat z_k))$ for $i=1..K$;\quad $v_k\gets\operatorname{mean}_i v^{(i)}_k$
\EndFor
\State \Return rollout $(\hat z,\hat r,q^{(\cdot)},q,v^{(\cdot)},v)$
\end{algorithmic}
\end{algorithm}

\begin{inwords}
The robot ``daydreams'': starting from what it really sees, it predicts the next hidden state, then the next, and
so on for $H$ steps. At every imagined step it also predicts the reward it would get and how good the situation is
(the value). During detection the actions are the ones the robot really executed, so the real outcome is known.
At run time they can be any candidate plan.
\end{inwords}

%==================================================================================================================
\section{The signals}\label{sec:signals}
%==================================================================================================================

All signals use the imagined rollout only: no simulator and no real future. Index $k=0,\dots,H-1$ (the learned
audit uses $k=0,\dots,L-1$).

\subsection{The core check: the Bellman residual}

For a consistent model, the value of a step equals its reward plus the discounted value of the next state:
\[
\Delta_k \;=\; \hat r_k+\gamma\,v_{k+1}-q_k \qquad(\text{signed residual, }\code{Rollout.signed\_residual}).
\]
If the imagination is consistent, $\Delta_k\approx0$ up to normal approximation error.

\begin{example}
The critic says the current step is worth $q_k=10$. The model imagines a reward $\hat r_k=1$ and says the next
imagined state is worth $v_{k+1}=3$. With $\gamma=0.95$: $\Delta_k=1+0.95\cdot3-10=-6.15$. The numbers do not add up,
so something in this imagined step is suspicious.
\end{example}

\subsection{Our 14 signals (\code{run\_preflight.py::plan\_signals}, \code{src/auditor/signals.py})}

\begin{center}
\renewcommand{\arraystretch}{1.35}
\begin{tabular}{@{}l>{$\displaystyle}l<{$}p{5.2cm}@{}}
\toprule
Name & \text{Formula} & What it catches\\
\midrule
A & \abs{\Delta_k} & one-step residual: sudden errors\\
A$n$ ($n{=}2,3,5,8$) & \abs{q_s-\Bigl(\sum_{j=s}^{k}\gamma^{j-s}\hat r_j+\gamma^{k-s+1}v_{k+1}\Bigr)},\ s=\max(0,k{-}n{+}1) & $n$-step residual ending at step $k$\\
Aa (anchored) & \abs{q_0-\Bigl(\sum_{j=0}^{k}\gamma^{j}\hat r_j+\gamma^{k+1}v_{k+1}\Bigr)} & imagined return vs the critic at the \emph{real} start: gradual drift\\
Ac (cumulative) & \abs{\sum_{j=0}^{k}\gamma^j\Delta_j} & signed sum: noise cancels, consistent drift adds up\\
Ao (optimistic) & \max(0,\Delta_k) & only over-promising steps (the ones a planner's argmax picks)\\
P & \sd_i\bigl(\hat r_k+\gamma v^{(i)}_{k+1}-q^{(i)}_k\bigr) & residual spread across critic heads\\
At & \abs{q_k-\Bigl(\hat r_k+\gamma\,\overline{\min}_{i<j}\bigl(v^{(i)}_{k+1,\text{tgt}},v^{(j)}_{k+1,\text{tgt}}\bigr)\Bigr)} & residual vs the critic's own training target (mean pairwise min of target heads)\\
B & \sd_i\,v^{(i)}_{k+1} & critic disagreement (also a baseline)\\
E & v_{k+1}+\beta\,\sd_i\,v^{(i)}_{k+1} & ELVIS-style optimistic value (also a baseline)\\
Ab & \abs{\Delta_k-b_k} & bias-corrected one-step residual\\
Acb & \abs{\sum_{j=0}^{k}\gamma^j(\Delta_j-b_j)} & bias-corrected cumulative residual\\
\bottomrule
\end{tabular}
\end{center}
That is 14 signals: A, A2, A3, A5, A8, Aa, Ac, Ao, P, At, B, E, Ab, Acb. Note A1 $=$ A, and Aa $=$ A at $k=0$.

\begin{inwords}
\textbf{A} checks one step at a time, so it catches sudden jumps. \textbf{Aa} always compares back to the real
starting point, where the critic is trustworthy, so small errors pile up against a fixed reference and slow drift
becomes visible. \textbf{Ac} adds the signed residuals: random noise has random signs and cancels, while a drift
pushes in one direction and grows. The others look at the same check from different angles (several steps, only
optimistic errors, disagreement between critic heads, the critic's own training target).
\end{inwords}

\begin{example}[title=Small example: why anchoring helps]
Suppose every imagined step is off by only $+0.3$ in value, below the alarm level of a one-step check. After 10
steps, A still sees $\approx0.3$ per step, but Aa compares the whole imagined return with the real start and sees
roughly $0.3\cdot(1+\gamma+\dots+\gamma^{9})\approx2.4$: the drift has become visible.
\end{example}

\subsection{Critic bias (\code{run\_preflight.py::calibration\_bias})}\label{sec:bias}

Even on correct steps, the mean critic sits slightly below $\hat r+\gamma v$, because TD-MPC2 trains toward a
\emph{minimum} of two heads. We measure this offset once per robot on verifiably clean calibration steps.

\begin{algorithm}[h]
\caption{\textsc{CriticBias}$(\theta,\mathcal C)$}\label{alg:bias}
\begin{algorithmic}[1]
\For{every length-$H$ window $w$ in the calibration episodes} \Comment{real start, real actions, real next observations}
  \State rollout $\gets\textsc{Imagine}(o^w_0,\mathbf a^w)$
  \State $e^w_k\gets\bigl\|\hat z_{k+1}-h_\theta(o^w_{k+1})\bigr\|$ for $k=0..H-1$ \Comment{latent error}
\EndFor
\State $\tau\gets P_{90}\bigl(\{e^w_0\}_w\bigr)$
\State clean$^w_k\gets\bigl[e^w_j\le\tau\ \ \forall j\le k\bigr]$ \Comment{clean prefix}
\State \Return $b_k=\med\{\Delta^w_k:\ \text{clean}^w_k\}$ for $k=0..H-1$
\end{algorithmic}
\end{algorithm}

\subsection{Baselines (same windows)}

\begin{center}
\renewcommand{\arraystretch}{1.35}
\begin{tabular}{@{}l>{$\displaystyle}l<{$}p{5.6cm}@{}}
\toprule
Name & \text{Formula} & Notes\\
\midrule
D & \tfrac1d\sum_{\text{dims}}\sd_m\, d^{(m)}_\theta(\hat z_k,a_k) & one-step dynamics-ensemble disagreement along the main rollout (MOPO / MACURA style); needs 4 extra dynamics heads\\
M & \sd_m\bigl(\hat r_k+\gamma V_\theta(d^{(m)}_\theta(\hat z_k,a_k))\bigr) & MOBILE-style spread of Bellman targets\\
B, E & \text{as above} & critic disagreement; ELVIS-style UCB (no extra models)\\
\bottomrule
\end{tabular}
\end{center}
The extra dynamics heads $d^{(1..4)}$ ($d^{(0)}=d_\theta$) are trained alongside the model on detached inputs, so
they never shape the encoder. Trajectory-level ensembles (Dt, Mt) are in Appendix~\ref{app:traj}.

%==================================================================================================================
\section{The answer key: what counts as a hallucination}\label{sec:labels}
%==================================================================================================================

Labels use only rewards (which the model is trained to predict) and the real rewards of the same actions from the
same real state (\code{run\_preflight.py::return\_error}). No probe, no reconstruction, no critic.

\begin{definition}[Imagination error]
\[
E_k=\abs{\sum_{j=1}^{k+1}\gamma^{\,j-1}\bigl(\hat r_j-r_j\bigr)},\qquad k=0,\dots,H-2,
\]
where $r_j$ is the real reward of action $a_j$ along the real trajectory. $\hat r_0$ is excluded because
$\hat z_0$ is the encoding of a real observation. Step $k$ judges the transition into $\hat z_{k+1}$.
\end{definition}

\begin{definition}[Thresholds and labels]
From the calibration half only, with $e_0$ the one-transition errors $E_0$ of all calibration windows:
$\varepsilon_{\text{hall}}=P_{99}(e_0)$, $\varepsilon_{\text{clean}}=P_{90}(e_0)$, and
\[
y_k=\begin{cases}1&E_k>\varepsilon_{\text{hall}}\quad\text{hallucinated}\\
0&E_k\le\varepsilon_{\text{clean}}\quad\text{correct}\\
\text{excluded}&\text{otherwise (ambiguous; never used for training or scoring).}\end{cases}
\]
\end{definition}

\begin{definition}[Sudden vs gradual]
With increments $I_j=E_j-E_{j-1}$ ($E_{-1}=0$) and $\kappa=P_{99}$ of calibration increments, a positive step $k$ is
\emph{sudden} if $I_j>\kappa$ for some $j\le k$, and \emph{gradual} otherwise.
\end{definition}

\begin{algorithm}[h]
\caption{\textsc{AnswerKey}(rollouts, episodes; thresholds)}\label{alg:labels}
\begin{algorithmic}[1]
\For{every window $w$} \State $E^w_k\gets\abs{\sum_{j=1}^{k+1}\gamma^{j-1}(\hat r^w_j-r^w_j)}$, $k=0..H-2$ \EndFor
\If{calibration} \State $\varepsilon_{\text{hall}}\gets P_{99}(\{E^w_0\})$;\ $\varepsilon_{\text{clean}}\gets P_{90}(\{E^w_0\})$ \EndIf
\State $y^w_k\gets 1$ if $E^w_k>\varepsilon_{\text{hall}}$;\ $0$ if $E^w_k\le\varepsilon_{\text{clean}}$;\ excluded otherwise
\State \Return labels $y$ (and thresholds)
\end{algorithmic}
\end{algorithm}

\begin{inwords}
A hallucination is an imagined future whose rewards add up to something really different from what actually
happened when the same actions were taken. ``Really different'' is measured against the model's own normal
one-step accuracy: only the worst 1\% of one-step errors counts as the bar for a hallucination, and anything
in the grey zone between ``normal'' and ``clearly wrong'' is left out so the test is not muddied.
\end{inwords}

\begin{example}
Calibration shows that one-step errors are usually tiny: $P_{90}=0.05$, $P_{99}=0.20$. A window whose imagined
reward sum is off by $0.6$ after 5 steps is labelled hallucinated at that step; one that is off by $0.03$ is
correct; one off by $0.10$ is ambiguous and skipped.
\end{example}

The simulator is deterministic given its state, so one real execution gives the real reward (a saved-state replay
reproduces a real step exactly on over 99\% of steps). In a stochastic environment, $r_j$ would be the average over
several replays.

%==================================================================================================================
\section{The learned audit: features and trees}\label{sec:features}
%==================================================================================================================

\subsection{Features: 31 numbers per imagined step (\code{run\_preflight.py::\_features})}

For the 14 signals $S=\{\text{A},\text{Aa},\text{Ac},\text{Ab},\text{Acb},\text{P},\text{B},\text{E},\text{At},
\text{Ao},\text{A2},\text{A3},\text{A5},\text{A8}\}$:
\[
x_k=\Bigl(\ \bigl[s_k,\ \max_{j\le k}s_j\bigr]_{s\in S},\ \ v_{k+1},\ \ \hat r_k,\ \ k\ \Bigr)\in\R^{14\cdot2+3}=\R^{31}.
\]
Each signal contributes its value now and the highest value so far in this rollout; then the imagined value, the
predicted reward and the step index.

\subsection{Fit once per robot, then freeze (\code{fit\_lba}, \code{make\_gbt})}\label{sec:trees}

\begin{algorithm}[h]
\caption{\textsc{FitTrees}}\label{alg:fit}
\begin{algorithmic}[1]
\State $X,y\gets[\ ],[\ ]$
\For{every length-$H$ window $w$ of the 25 calibration episodes}
  \State rollout $\gets\textsc{Imagine}(o^w_0,\mathbf a^w)$;\ signals with bias $b$;\ $E^w\gets$ imagination error
  \For{$k=0$ to $L-1$}
    \If{$y^w_k$ is not excluded} \State append $x^w_k$ to $X$ and $y^w_k$ to $y$ \EndIf
  \EndFor
\EndFor
\State \Return gradient-boosted trees fitted on $(X,y)$ \Comment{200 trees, 15 leaves, learning rate 0.05, $\ell_2=1$, balanced classes, seed 0}
\end{algorithmic}
\end{algorithm}

The tree settings are fixed for every robot and were never tuned on test data. The trees output
$p_k=\Pr(y_k=1\mid x_k)$.

\subsection{Predict (no labels, no simulator)}

\begin{algorithm}[h]
\caption{\textsc{Detect}$(f,b,o_t,\mathbf a)$}\label{alg:detect}
\begin{algorithmic}[1]
\State rollout $\gets\textsc{Imagine}(o_t,\mathbf a)$;\ compute the 14 signals with bias $b$
\State \Return $p_k=f(x_k)$ for $k=0..L-1$ \Comment{probability that imagined step $k$ is hallucinated}
\end{algorithmic}
\end{algorithm}

\begin{inwords}
No single signal is right everywhere: one is good at sudden jumps, another at slow drift, another near the end of a
rollout. The small decision trees learn, from the calibration episodes, which combinations really mean
``hallucination'' for this robot. After that they are frozen and only read the signals.
\end{inwords}

\paragraph{Variants reported.} The same features with logistic regression (``linear''); single raw signals such as
Aa; and controls: \emph{NoBellman} trees with only $(v_{k+1},\hat r_k,k)$, and \emph{BellmanOnly} trees with the
signals and $k$ only. The Bellman signals add about $+0.05$ AUROC over NoBellman on both benchmarks.

%==================================================================================================================
\section{Evaluation}\label{sec:eval}
%==================================================================================================================

\subsection{Protocol per trained robot (\code{run\_preflight.py::evaluate\_model})}

Record 50 episodes with the agent's own planner (evaluation mode, seeds 2000--2049). Fit everything (bias,
thresholds, trees) on the first 25; nothing from the test half is used for fitting. For every length-$H$ window of
every test episode (start at each real step, real actions): imagine, label with the answer key, and score with
every detector (ours and baselines).

\subsection{Metric: step-stratified AUROC (\code{src/preflight/metrics.py})}

For scores $s$ and labels $y$, with positives $\mathcal P$ and negatives $\mathcal N$:
\[
\operatorname{AUROC}=\frac{1}{|\mathcal P||\mathcal N|}\sum_{p\in\mathcal P}\sum_{n\in\mathcal N}
\Bigl(\mathbb 1[s_p>s_n]+\tfrac12\,\mathbb 1[s_p=s_n]\Bigr).
\]
Errors naturally grow with the imagined step, so a detector that only knows ``later is worse'' would look good.
We therefore compare only within the same step $k$ and pool:
\[
\operatorname{AUROC}_{\text{strat}}=\frac{\sum_k n_k\operatorname{AUROC}_k}{\sum_k n_k},\qquad n_k=|\mathcal P_k|\,|\mathcal N_k|.
\]

\begin{inwords}
AUROC answers: if I pick one hallucinated step and one correct step at random, how often does the detector give the
hallucinated one the higher alarm? $0.5$ is a coin flip, $1.0$ is perfect. ``Stratified'' means we only ever compare
steps at the same distance into the future, so the detector cannot win just by knowing that later steps are worse.
\end{inwords}

\begin{example}
3 hallucinated steps with scores $0.9,0.7,0.4$ and 2 correct steps with scores $0.5,0.2$. Of the $3\times2=6$ pairs,
the hallucinated one scores higher in 5 ($0.4<0.5$ is the only miss), so AUROC $=5/6\approx0.83$.
\end{example}

\subsection{Uncertainty and comparisons}

\begin{itemize}[nosep]
\item \textbf{Bootstrap CI:} resample whole test episodes with replacement 1000 times; the 95\% interval is
$[P_{2.5},P_{97.5}]$ of $\operatorname{AUROC}_{\text{strat}}$ over resamples (episodes, not steps, because steps of one
episode are correlated).
\item \textbf{Paired difference:} the same resampled episodes for both detectors; ours is significantly better
if the 95\% interval of $\operatorname{AUROC}_{\text{ours}}-\operatorname{AUROC}_{\text{rival}}$ excludes 0.
\item Reported for all positives, sudden-only positives and gradual-only positives.
\end{itemize}

\subsection{Robots and confirmation}

Main tables: 4 ManiSkill tasks (PushCube, StackCube, PickSingleYCB, PegInsertionSide), seeds 10--40, grounded and
standard critics. Confirmation: new seed-50 robots with every rule fixed in advance
(\code{docs/preregistration\_seed50.md}): all criteria passed (ours 0.81 vs best rival 0.73). Second benchmark: the
released TD-MPC2 multi-task models on all 30 DeepMind Control tasks (\code{src/dmc\_audit.py}).

%==================================================================================================================
\section{Run time}\label{sec:runtime}
%==================================================================================================================

At run time the detector needs only the trained world model, the bias $b$ and the trees $f$: Algorithm~\ref{alg:detect}
on any candidate action sequence. No simulator, no real future and no rewards are used.

For many candidates at once (e.g.\ all 512 MPPI candidates per planning iteration), \code{src/planning/gpu\_audit.py}
evaluates the same trees exactly on the GPU; every job checks the GPU copy against scikit-learn before use (maximum
difference $0.0$). Uses of the detector beyond detection (planning, training, shift monitoring, policy evaluation,
per-signal audits) were tested separately; their designs and results are in \code{docs/preregistration\_*.md}.

%==================================================================================================================
\appendix
\section{Extensions (October 2026)}
%==================================================================================================================

\subsection{Compounding error along long rollouts (\code{src/tools/compounding.py}, \code{dmc\_compounding.py})}

Rollouts of $H=24$ and $40$ steps, same labels. For each signal $s$ and robot, the mean signal on hallucinated
windows at step $k$ is $\bar s_k=\operatorname{mean}\{s_k: y_k=1\}$, and its growth from the first to the last 3 steps is
\[
G_s=\frac{\frac13\sum_{k\in\text{last 3}}\bar s_k}{\frac13\sum_{k\in\text{first 3}}\bar s_k},
\qquad\text{compared with the growth of the true error }\ G_E=\frac{\operatorname{mean}_{\text{last 3}}E_k}{\operatorname{mean}_{\text{first 3}}E_k}.
\]
We report the mean of $G_s$ over robots (or tasks) with a 95\% bootstrap interval over robots (2000 resamples), the
number of robots with $G_s>1$, and how often $G_{\text{ours}}>G_{\text{ensemble}}$. Per-step stratified AUROC
measures ranking separately from loudness.

\begin{inwords}
As the robot imagines further, its errors pile up (in our data 18 to 112 times larger at the end than at the
start). A good alarm should get louder as the error grows. $G_s$ measures how much louder each alarm gets.
\end{inwords}

\subsection{Trajectory-level ensembles and shared bias (\code{signals.py::member\_rollouts})}\label{app:traj}

Each dynamics head imagines the whole rollout on its own (PETS-style trajectory sampling):
$\hat z^{(m)}_0=\hat z_0$, $\hat z^{(m)}_{k+1}=d^{(m)}_\theta(\hat z^{(m)}_k,a_k)$, $m=0..4$.
\begin{align*}
\text{Dt}_k&=\tfrac1d\textstyle\sum_{\text{dims}}\sd_m\,\hat z^{(m)}_{k+1} &&\text{(spread of the imagined states)}\\
G^{(m)}_k&=\textstyle\sum_{j\le k}\gamma^j R_\theta(\hat z^{(m)}_j,a_j)+\gamma^{k+1}V_\theta(\hat z^{(m)}_{k+1}),\qquad
\text{Mt}_k=\sd_m\,G^{(m)}_k &&\text{(STEVE-style spread of the heads' own returns)}
\end{align*}
\textbf{Shared-bias diagnostic} (\code{run\_preflight.py::shared\_bias}): with member errors
$e^{(m)}_k=\hat z^{(m)}_{k+1}-h_\theta(o_{k+1})$ and their mean $\bar e_k$, we log the error of the ensemble mean
$\|\bar e_k\|$, the mean member error, the spread, and the same at the level of reward sums.

\subsection{Why the anchored check must grow, and when ensembles cannot see an error}

Let $G_k(x)=\sum_{j\le k}\gamma^jR_\theta(x_j,a_j)+\gamma^{k+1}V_\theta(x_{k+1})$ along a latent path $x$; then
$\text{Aa}_k=\abs{\Qbar(\hat z_0,a_0)-G_k(\hat z)}$.

\begin{proposition}[Anchored residual $=$ imagination error $\pm$ a dynamics-free floor]
Define the imagination error in value units $I_k=G_k(\hat z)-G_k(z)$ (zero if the imagined latents equal the real
encoded ones $z_j=h_\theta(o_j)$) and the critic floor $C_k=\Qbar(\hat z_0,a_0)-G_k(z)$, which does not involve the
dynamics at all. Then $\Qbar(\hat z_0,a_0)-G_k(\hat z)=C_k-I_k$ exactly, hence
\[
\bigl|\,|I_k|-|C_k|\,\bigr|\ \le\ \text{Aa}_k\ \le\ |I_k|+|C_k|.
\]
\end{proposition}
\noindent So once the value-relevant drift $|I_k|$ exceeds the floor, Aa must grow with it. A grounded critic keeps
$|C_k|$ small. $|C_k|$ is largest where the value changes sharply (contacts, value cliffs), which is where false
alarms concentrate.

\begin{proposition}[Ensemble spread cannot see shared error]
For $K$ members with errors $e^{(m)}$ and mean $\bar e$:
$\ \frac1K\sum_m\|e^{(m)}\|^2=\|\bar e\|^2+\frac1K\sum_m\|e^{(m)}-\bar e\|^2$.
Any alarm that depends only on disagreement is unchanged if all members shift by the same offset.
\end{proposition}
\noindent Whether shared error dominates in practice is measured by the shared-bias diagnostic above (result
pending at the time of writing).

\subsection{Hallucination-aware value expansion (\code{src/training/value\_expansion.py})}

For a real transition with reward $r$ and real next latent $z'$, imagine $h$ steps with the policy:
$\hat z_0=z'$, $\hat z_{k+1}=d_\theta(\hat z_k,\pi(\hat z_k))$, $\hat r_k=R_\theta(\hat z_k,\pi(\hat z_k))$. The $j$-step
targets are
\[
G_j=r+\gamma\Bigl(\sum_{k<j}\gamma^k\hat r_k+\gamma^j\,Q_{\bar\theta}\bigl(\hat z_j,\pi(\hat z_j)\bigr)\Bigr),\qquad j=0..h,
\]
($G_0$ is the standard TD-MPC2 target). Each imagined step $k$ gets a trust $t_k\in[0,1]$; imagination is cut at
the first untrusted step, and the critic's target is the expected return under that cut:
\[
T_0=1,\quad T_j=\prod_{k<j}t_k,\qquad
G=\sum_{j<h}T_j(1-t_j)\,G_j+T_h\,G_h .
\]
Trust: \emph{none} uses $G_0$; \emph{fixed} $t\equiv1$ (uses $G_h$); \emph{const} $t\equiv0.5$; \emph{D} and \emph{lba}
map a raw score (ensemble disagreement, or our audit's probability) through the same per-step empirical CDF over
recent scores, $t_k=1-\widehat{F}_k(\text{score}_k)$, so both have the same average trust and differ only in
\emph{where} they trust.

\begin{inwords}
Instead of learning only from the next real step, the critic also learns from a few imagined steps ahead, which is
faster. But each imagined step is trusted only as much as the detector allows: a suspicious step cuts the
imagination short. The pilot showed that trusting imagination fully is a disaster (final success 0.22), while
careful trust helps; the full study (55 runs, pre-registered) is running.
\end{inwords}

\end{document}
